\subsection{Постановка научной задачи и концептуальная модель SmartSQL}
\subsubsection{Результаты анализа и выявленные ограничения}

Проведённый обзор позволяет выделить несколько областей, которые в рассмотренных исследованиях чаще анализируются по отдельности. Первая область — формализация пользовательского намерения и правильное сопоставление естественно-языкового запроса со схемой. Вторая — адаптация интерфейса к уровню подготовки и характеру задачи. Третья — формальная проверка допустимости и последствия исполнения операций. Четвёртая — переносимость подхода между различными источниками структурированных данных.

Работы по естественно-языковому доступу демонстрируют разнообразие способов формализации запроса \cite{posevkin2018,almuhammed2022,majeed2022}. Исследования адаптации рассматривают модель пользователя и управление представлением функций \cite{tagirova2022,abrahao2021}. Публикации по работе с данными изучают интеграцию, обработку и перенос между СУБД \cite{zayakin2022,voronchikhin2024}. Обобщение этих направлений позволяет сформулировать исследовательскую проблему: \emph{как согласовать доступность естественно-языкового взаимодействия, вариативность пользовательского интерфейса и проверяемое управление операциями над разными структурами данных в рамках единой концептуальной модели?}

Нельзя заранее предполагать, что объединение перечисленных функций улучшит все показатели одновременно. Например, дополнительный этап проверки способен повысить контролируемость изменений, но увеличить время выполнения задачи; автоматическое переключение режима может сократить число видимых элементов, но дезориентировать часть пользователей. Поэтому научная постановка требует явного определения компромиссов и критериев сравнения.

С.~Н.~Ушаков и А.~О.~Савельев в сравнительном обзоре автоматизированного извлечения знаний из научных текстов демонстрируют важность этапов предварительной обработки, выделения объектов и представления результатов \cite{ushakov2024}. Хотя предмет их исследования отличается от управления реляционными БД, сама идея стадийного конвейера полезна как методологическая аналогия: каждый этап должен иметь собственный вход, выход и показатели качества. Эта аналогия не подменяет доказательства применимости конкретного алгоритма к SQL.

\subsubsection{Предлагаемая концептуальная схема}

На основании обзора предлагается рассмотреть последовательность из пяти логических компонентов:
\begin{enumerate}
\item \textbf{Представление данных и метаданных.} Описывает источники, схемы, ключи, типы полей и доступные операции, не отождествляя все СУБД с одним SQL-диалектом.
\item \textbf{Интерпретация намерения.} Принимает формулировку пользователя или структурированное действие графического интерфейса и формирует внутреннее представление операции.
\item \textbf{Проверка и уточнение.} Выявляет неопределённости, сопоставляет операцию с ограничениями доступа и определяет необходимость подтверждения.
\item \textbf{Исполнение.} Передаёт разрешённую операцию конкретному провайдеру данных и получает результат либо диагностическое сообщение.
\item \textbf{Адаптивное представление.} Отображает результаты и объяснения в форме, соответствующей выбранному сценарию и уровню пользователя.
\end{enumerate}

Под \emph{внутренним представлением} здесь понимается не обязательный язык программирования, а логический объект, содержащий цель операции, упоминания таблиц и полей, условия отбора, сортировку, ограничения и тип воздействия на данные. Такой объект может быть проверен до преобразования в SQL конкретного провайдера. Принципиальная особенность предлагаемой схемы состоит в том, что и графический, и естественно-языковой интерфейсы направляют запросы на общую проверку. При этом алгоритм интерпретации не получает прямого доступа к исполнению произвольных команд.

\subsubsection{Исследовательские гипотезы}

Для следующего этапа научно-исследовательской работы предлагается проверить следующие гипотезы:
\begin{enumerate}
\item \textbf{Гипотеза H1.} В типовых сценариях чтения и фильтрации таблиц естественно-языковой ввод с визуализацией предполагаемой операции может повысить результативность начинающих пользователей по сравнению с самостоятельным написанием SQL.
\item \textbf{Гипотеза H2.} Интерфейс с управляемыми уровнями сложности может уменьшить время поиска нужной функции для начинающих, не приводя к ухудшению результативности опытных пользователей в профессиональных задачах.
\item \textbf{Гипотеза H3.} Разделение генерации и проверки операций уменьшает число исполненных недопустимых действий по сравнению с непосредственным выполнением машинно сформированных запросов, при возможном увеличении времени взаимодействия.
\item \textbf{Гипотеза H4.} Использование метаданных схемы и общего внутреннего представления может уменьшить долю ошибок при переносе однотипных пользовательских сценариев между поддерживаемыми СУБД.
\end{enumerate}

Каждая гипотеза проверяема и допускает отрицательный результат. В частности, H1 требует оценки не только скорости, но и правильности возвращаемых данных; H3 — перечня запрещённых и разрешённых сценариев; H4 — заранее зафиксированного набора поддерживаемых операций и честного учёта несовместимых возможностей разных систем.

\subsubsection{Предварительный протокол проверки}

Для будущего эксперимента следует сформировать набор тестовых баз с документированными схемами, а также корпус заданий разных уровней сложности. Сценарии должны включать однозначные запросы, намеренно неоднозначные формулировки и потенциально опасные действия. Для каждого задания необходим эталонный смысловой результат, который не сводится только к строковому совпадению SQL. Эквивалентные запросы могут отличаться синтаксисом, а одинаково выглядящие запросы — иметь разные последствия на различных данных.

Предлагается три группы критериев. \emph{Семантическое качество}: доля правильно интерпретированных запросов, корректность выбранных таблиц и атрибутов, частота необходимых уточнений. \emph{Управляемость и безопасность}: доля заблокированных недопустимых операций, число ложных блокировок, ясность подтверждения. \emph{Удобство}: время решения, число ошибок, субъективная нагрузка и удовлетворённость. Необходимо фиксировать не только средние значения, но и разброс результатов, поскольку агрегированные показатели могут скрывать существенные различия между группами пользователей.

В исследовании не планируется считать «универсальность» бинарным свойством. Для каждого подключаемого источника данных должна быть составлена матрица поддерживаемых возможностей: чтение схемы, выборка, фильтрация, изменения строк, транзакции, изменение структуры и экспорт. Такое представление позволит ограничить выводы реально исследованными операциями и избежать необоснованных утверждений о полной переносимости.

\subsubsection{Границы исследования первого семестра}

Настоящий отчёт фиксирует теоретическую основу, корпус литературы, концептуальную схему и план оценки. Он не содержит эмпирических данных об эффективности SmartSQL, не утверждает достижения заявленных показателей и не претендует на доказанную научную новизну. До проведения эксперимента обоснованными являются только постановка проблемы, синтез рассмотренных подходов и формулирование проверяемых предположений.

Дальнейшая работа должна включать уточнение границ автоматизации, подготовку корпуса запросов на русском языке, формальное описание политик допуска операций, разработку прототипов интерфейса и сопоставимое испытание альтернативных сценариев. Итоги этих действий станут основой для последующего отчёта, тезисов конференции и научной публикации при условии получения воспроизводимых результатов.